% 子文件：10.向量场与第二型积分.tex
% 主文档使用 \input{10.向量场与第二型积分} 引入，不含 document 环境

\section{向量场}
\subsection{向量场}

\subsection{向径的导数}

\section{第二型曲线积分}
\subsection{变力做功与第二型曲线积分的概念}

\subsection{第二型曲线积分的计算}

\subsection{第二型曲线积分与第一型曲线积分的关系}

\section{格林公式、平面流速场的环量与旋度}
\subsection{格林公式}

\subsection{平面流速场的环量与旋度}

\section{平面曲线积分与路径无关的条件、保守场}
\subsection{平面曲线积分与路径无关的条件}

定理10.2\quad

\subsection{保守场、原函数、全微分方程}

\section{第二型曲面积分}
\subsection{预备知识}
\begin{enumerate}
    \item 有向曲面
    \item 有向曲面在坐标面上的投影
\end{enumerate}

\subsection{第二型曲面积分的概念}

\subsection{第二型曲面积分的计算}

\section{高斯公式、通量与散度}
\subsection{高斯公式}

定理10.3\quad 则有高斯公式
\[
\oiint_\Sigma P\,dydz+Q\,dzdx+R\,dxdy
=\iiint_V\left(\frac{\partial P}{\partial x}+\frac{\partial Q}{\partial y}+\frac{\partial R}{\partial z}\right)dV
\]
即
\[
\oiint_\Sigma \boldsymbol{F}\cdot d\boldsymbol{S}
=\iiint_V \mathrm{div}\,\boldsymbol{F}\,dV
\]
其中

\subsection{向量场的通量与散度}

定义10.4\quad

\section{斯托克斯公式、环量与旋度}
\subsection{斯托克斯公式}

\subsection{向量场的环量与旋度}